%% Abstract
%%
%% Traduction anglaise fidèle et de qualité du résumé de la recherche écrit en français et non une traduction littérale. 
%%

\chapter*{ABSTRACT}\thispagestyle{headings}
\addcontentsline{toc}{compteur}{ABSTRACT}
%
\begin{otherlanguage}{english}

Wireless Sensor Networks (WSNs) deployed in disaster-prone environments provide critical real-time monitoring of natural phenomena such as earthquakes, floods, and landslides. However, three fundamental challenges hinder their large-scale deployment: (i)~limited network lifetime imposed by non-rechargeable batteries, (ii)~the full lifecycle carbon footprint of massively distributed devices, and (iii)~the vulnerability of AI-learned routing policies to adversarial attacks and model hallucinations. This thesis, organized as a thesis by articles, proposes three complementary frameworks that sequentially address these challenges in the context of AI-native 6G sensing networks.

The first article introduces GAPF (Graph-Attentive Policy Fusion), a lightweight meta-learning framework (1,473 parameters) that fuses frozen multi-agent reinforcement learning expert policies (QMIX, QTRAN) through a 2-layer Graph Convolutional Network encoder and a Compact Attention Selector controller. The bi-level architecture combines a Monotonic Q$\cdot$K Attention Network for inner-loop reward scalarization with Evolution Strategies for outer-loop optimization. Experimental results demonstrate a packet delivery ratio (PDR) $\geq 98\%$, 2--3.5$\times$ lower energy consumption, and a 13$\times$ smaller memory footprint compared to baselines, while maintaining near real-time inference.

The second article presents CALASH (Carbon-Aware Lifecycle-Autonomous Self-Healing), the first WSN routing protocol that jointly minimizes Lifecycle Carbon Intensity (LCI)---defined as the embodied, operational, and end-of-life CO$_2$ per useful delivered packet---through four integrated pillars: (i)~CADR, which modulates compressive-sensing ratios based on grid carbon intensity; (ii)~CARE, a hybrid Lyapunov--DQN router with provable carbon-budget compliance guarantees; (iii)~SHDR, an autonomic MAPE-K self-healing loop; and (iv)~LSE, an ISO~14040-compliant lifecycle sustainability engine. Operating over a dual-band 6G physical layer (sub-THz / sub-6\,GHz), CALASH achieves 60\% longer lifetime than LEACH (1,486 vs.\ 931 rounds), 82.9\% PDR, and 33\% lower LCI (8.99 vs.\ 13.42\,gCO$_2$eq/pkt).

The third article proposes CASTER-ZT (Zero-Trust Shielded Graph-Structured Decision Control), the first zero-trust decision control framework for disaster-monitoring sensing networks. CASTER-ZT couples a 2-layer GraphSAGE encoder, a learned recovery policy, a neural dual-hypothesis trust-risk evaluator, MC-Dropout epistemic uncertainty quantification, and a deterministic shield producing graduated decisions (allow, scope-reduce, defer, escalate, block). Ten formal propositions establish monotone conservatism, unauthorized-actuation exclusion, and conformal coverage guarantees. A 2,420-run campaign demonstrates 74.2--98.3\% rogue-policy detection with zero false positives, a weighted recovery score $\omega_{\text{rec}} = 0.733$ (3.7--5.2$\times$ higher than baselines), and 3.07\,ms per-decision latency within O-RAN's 10\,ms near-RT budget.

Collectively, the three articles establish a complete stack---energy efficiency, carbon sustainability, and zero-trust security---for intelligent routing in 6G-integrated disaster-monitoring sensor networks.

\end{otherlanguage}
